% GENERATED FILE. Edit canonical lessons, then run npm run generate:books.
% canonical-source-hash: fnv1a64:eb943f9d1f4aa876
% canonical-lessons: TE-C427-pattubattu, TE-C427-kudurcu, TE-C427-samadhanaparacu, TE-C427-santrptiparacu, TE-C427-madhyavarti

\chapter{Coming to Terms}
\label{ch:te-coming-to-terms}
\hlchaptermodality{427}

\begin{chapteropening}
\textbf{By the end of this chapter:} \emph{I can insist, arrange, reassure and satisfy.}

Complete the last lesson of chapter 427: I can insist, arrange, reassure and satisfy.
\end{chapteropening}

\section[\texorpdfstring{\te{పట్టుబట్టు} (paṭṭubaṭṭu)}{పట్టుబట్టు (paṭṭubaṭṭu)}]{\te{పట్టుబట్టు} (paṭṭubaṭṭu) — to insist}

\label{lesson:TE-C427-pattubattu}



\begin{quote}
Before the new one: say the Telugu for deliberately, then the Telugu for a beginning.
\end{quote}

\subsection*{You'll want to know: \te{పట్టుబట్టు}}
\textbf{\te{పట్టుబట్టు}} — \emph{paṭṭubaṭṭu} — \textquotedblleft{}to insist\textquotedblright{}.

\begin{quote}
\te{ఆమె} \te{ఆ} \te{రోజే} \te{వెళ్ళాలని} \te{పట్టుబట్టింది}.
\end{quote}

\begin{quote}
\emph{āme ā rōjē veḷḷālani paṭṭubaṭṭindi.}
\end{quote}

\begin{quote}
She insisted on going that very day.
\end{quote}

Past (I …): \te{పట్టుబట్టాను}, \emph{paṭṭubaṭṭānu}.

\te{పట్టు}, a grip, + \te{పట్టు}, to hold: to keep a firm grip on one’s point.

\begin{cousinweb}
\textbf{\te{పట్టుబట్టు}} (\emph{paṭṭubaṭṭu}) is native Dravidian, a compound of \emph{paṭṭu}, \textquotedblleft{}a hold, a grasp; perseverance\textquotedblright{}, and \emph{paṭṭu}, \textquotedblleft{}to hold, to seize\textquotedblright{} (DEDR 4034).
\end{cousinweb}

\begin{grammarlens}[title={What you've built}]
One more word for saying what happens if.
\end{grammarlens}

\subsection*{Your turn}
\begin{itemize}
\raggedright
  \item \emph{Say it:} \emph{paṭṭubaṭṭu}
  \item \emph{Say it:} \emph{paṭṭubaṭṭu}, once more
  \item \emph{Say it:} \emph{āme ā rōjē veḷḷālani paṭṭubaṭṭindi.} — She insisted on going that very day.
  \item \emph{Recall:} say the Telugu for deliberately, then the Telugu for a beginning, then say \emph{paṭṭubaṭṭu} again
\end{itemize}

\begin{tcolorbox}[breakable,colback=teal!4,colframe=teal!35!black,title={Before you move on}]
What does \te{పట్టుబట్టు} mean? (\textquotedblleft{}To insist\textquotedblright{}.) Say it once more.
\end{tcolorbox}
\section[\texorpdfstring{\te{కుదుర్చు} (kudurcu)}{కుదుర్చు (kudurcu)}]{\te{కుదుర్చు} (kudurcu) — to fix up; to arrange (a deal)}

\label{lesson:TE-C427-kudurcu}



\begin{quote}
Before the new one: say the Telugu for a beginning, then the Telugu for to insist.
\end{quote}

\subsection*{You'll want to know: \te{కుదుర్చు}}
\textbf{\te{కుదుర్చు}} — \emph{kudurcu} — \textquotedblleft{}to fix up; to arrange (a deal)\textquotedblright{}.

\begin{quote}
\te{మా} \te{మామయ్య} \te{ఆ} \te{పెళ్ళి} \te{సంబంధం} \te{కుదిర్చాడు}.
\end{quote}

\begin{quote}
\emph{mā māmayya ā peḷḷi sambandhaṁ kudircāḍu.}
\end{quote}

\begin{quote}
My uncle arranged that match.
\end{quote}

Past (I …): \te{కుదిర్చాను}, \emph{kudircānu}.

Also for deals: \te{ఒప్పందం} \te{కుదుర్చు}, to broker an agreement.

\begin{cousinweb}
\textbf{\te{కుదుర్చు}} (\emph{kudurcu}) is native Dravidian: DEDR 1709 lists Telugu \emph{kudurcu} beside \emph{kuduru}, \textquotedblleft{}to become settled, to be adjusted or arranged\textquotedblright{}, with Tamil \emph{kutir}, \textquotedblleft{}to be settled\textquotedblright{}.
\end{cousinweb}

\begin{grammarlens}[title={What you've built}]
One more word for saying what happens if.
\end{grammarlens}

\subsection*{Your turn}
\begin{itemize}
\raggedright
  \item \emph{Say it:} \emph{kudurcu}
  \item \emph{Say it:} \emph{kudurcu}, once more
  \item \emph{Say it:} \emph{mā māmayya ā peḷḷi sambandhaṁ kudircāḍu.} — My uncle arranged that match.
  \item \emph{Recall:} say the Telugu for a beginning, then the Telugu for to insist, then say \emph{kudurcu} again
\end{itemize}

\begin{tcolorbox}[breakable,colback=teal!4,colframe=teal!35!black,title={Before you move on}]
What does \te{కుదుర్చు} mean? (\textquotedblleft{}To fix up; to arrange (a deal)\textquotedblright{}.) Say it once more.
\end{tcolorbox}
\section[\texorpdfstring{\te{సమాధానపరచు} (samādhānaparacu)}{సమాధానపరచు (samādhānaparacu)}]{\te{సమాధానపరచు} (samādhānaparacu) — to reassure; to calm down}

\label{lesson:TE-C427-samadhanaparacu}



\begin{quote}
Before the new one: say the Telugu for to insist, then the Telugu for to fix up.
\end{quote}

\subsection*{You'll want to know: \te{సమాధానపరచు}}
\textbf{\te{సమాధానపరచు}} — \emph{samādhānaparacu} — \textquotedblleft{}to reassure; to calm down\textquotedblright{}.

\begin{quote}
\te{ఏడుస్తున్న} \te{పిల్లవాడిని} \te{అమ్మ} \te{సమాధానపరిచింది}.
\end{quote}

\begin{quote}
\emph{ēḍustunna pillavāḍini amma samādhānaparicindi.}
\end{quote}

\begin{quote}
Mother calmed the crying child.
\end{quote}

Past (I …): \te{సమాధానపరిచాను}, \emph{samādhānaparicānu}.

\te{సమాధానం}, an answer, and also calm, + -\te{పరచు}, an ending that makes a verb.

\begin{cousinweb}
\textbf{\te{సమాధానపరచు}} (\emph{samādhānaparacu}) is built from Sanskrit \emph{samādhāna}, \textquotedblleft{}putting together; reconciliation\textquotedblright{} (CDIAL 13210), and -\te{పరచు}, a native ending that makes a verb (DEDR 3853 lists \emph{paṟacu}, \textquotedblleft{}to cause to feel\textquotedblright{}, beside \emph{paḍu}).
\end{cousinweb}

\begin{grammarlens}[title={What you've built}]
One more word for saying what happens if.
\end{grammarlens}

\subsection*{Your turn}
\begin{itemize}
\raggedright
  \item \emph{Say it:} \emph{samādhānaparacu}
  \item \emph{Say it:} \emph{samādhānaparacu}, once more
  \item \emph{Say it:} \emph{ēḍustunna pillavāḍini amma samādhānaparicindi.} — Mother calmed the crying child.
  \item \emph{Recall:} say the Telugu for to insist, then the Telugu for to fix up, then say \emph{samādhānaparacu} again
\end{itemize}

\begin{tcolorbox}[breakable,colback=teal!4,colframe=teal!35!black,title={Before you move on}]
What does \te{సమాధానపరచు} mean? (\textquotedblleft{}To reassure; to calm down\textquotedblright{}.) Say it once more.
\end{tcolorbox}
\section[\texorpdfstring{\te{సంతృప్తిపరచు} (santṛptiparacu)}{సంతృప్తిపరచు (santṛptiparacu)}]{\te{సంతృప్తిపరచు} (santṛptiparacu) — to satisfy; to please}

\label{lesson:TE-C427-santrptiparacu}



\begin{quote}
Before the new one: say the Telugu for to fix up, then the Telugu for to reassure.
\end{quote}

\subsection*{You'll want to know: \te{సంతృప్తిపరచు}}
\textbf{\te{సంతృప్తిపరచు}} — \emph{santṛptiparacu} — \textquotedblleft{}to satisfy; to please\textquotedblright{}.

\begin{quote}
\te{అందరినీ} \te{సంతృప్తిపరచడం} \te{కష్టం}.
\end{quote}

\begin{quote}
\emph{andarinī santṛptiparacaḍaṁ kaṣṭaṁ.}
\end{quote}

\begin{quote}
It is hard to satisfy everyone.
\end{quote}

Past (I …): \te{సంతృప్తిపరిచాను}, \emph{santṛptiparicānu}.

\te{సంతృప్తి}, satisfaction, + -\te{పరచు}. \te{నాకు} \te{సంతృప్తిగా} \te{ఉంది}, I am satisfied.

\begin{cousinweb}
\textbf{\te{సంతృప్తిపరచు}} (\emph{santṛptiparacu}) is built from Sanskrit \emph{saṁ-}, \textquotedblleft{}wholly\textquotedblright{}, and \emph{tṛpti}, \textquotedblleft{}satisfaction, contentment\textquotedblright{} (Monier-Williams), with -\te{పరచు}, a native ending that makes a verb (DEDR 3853 lists \emph{paṟacu}, \textquotedblleft{}to cause to feel\textquotedblright{}, beside \emph{paḍu}).
\end{cousinweb}

\begin{grammarlens}[title={What you've built}]
One more word for saying what happens if.
\end{grammarlens}

\subsection*{Your turn}
\begin{itemize}
\raggedright
  \item \emph{Say it:} \emph{santṛptiparacu}
  \item \emph{Say it:} \emph{santṛptiparacu}, once more
  \item \emph{Say it:} \emph{andarinī santṛptiparacaḍaṁ kaṣṭaṁ.} — It is hard to satisfy everyone.
  \item \emph{Recall:} say the Telugu for to fix up, then the Telugu for to reassure, then say \emph{santṛptiparacu} again
\end{itemize}

\begin{tcolorbox}[breakable,colback=teal!4,colframe=teal!35!black,title={Before you move on}]
What does \te{సంతృప్తిపరచు} mean? (\textquotedblleft{}To satisfy; to please\textquotedblright{}.) Say it once more.
\end{tcolorbox}
\section[\texorpdfstring{\te{మధ్యవర్తి} (madhyavarti)}{మధ్యవర్తి (madhyavarti)}]{\te{మధ్యవర్తి} (madhyavarti) — a mediator; a go-between}

\label{lesson:TE-C427-madhyavarti}



\begin{quote}
Before the new one: say the Telugu for to reassure, then the Telugu for to satisfy.
\end{quote}

\subsection*{You'll want to know: \te{మధ్యవర్తి}}
\textbf{\te{మధ్యవర్తి}} — \emph{madhyavarti} — \textquotedblleft{}a mediator; a go-between\textquotedblright{}.

\begin{quote}
\te{రెండు} \te{కుటుంబాల} \te{మధ్య} \te{ఒక} \te{మధ్యవర్తి} \te{మాట్లాడాడు}.
\end{quote}

\begin{quote}
\emph{reṇḍu kuṭumbāla madhya oka madhyavarti māṭlāḍāḍu.}
\end{quote}

\begin{quote}
A mediator spoke for both families.
\end{quote}

Plural: \te{మధ్యవర్తులు}, \emph{madhyavartulu}.

\te{మధ్య}, between, + \te{వర్తి}, one who stays.

\begin{cousinweb}
\textbf{\te{మధ్యవర్తి}} (\emph{madhyavarti}) is built from Sanskrit \emph{madhya}, \textquotedblleft{}middle\textquotedblright{} (CDIAL 9804), and \emph{vartin}, \textquotedblleft{}abiding, situated in\textquotedblright{} (Monier-Williams): one who stands in the middle.
\end{cousinweb}

\begin{grammarlens}[title={What you've built}]
One more word for saying what happens if.
\end{grammarlens}

\subsection*{Your turn}
\begin{itemize}
\raggedright
  \item \emph{Say it:} \emph{madhyavarti}
  \item \emph{Say it:} \emph{madhyavarti}, once more
  \item \emph{Say it:} \emph{reṇḍu kuṭumbāla madhya oka madhyavarti māṭlāḍāḍu.} — A mediator spoke for both families.
  \item \emph{Recall:} say the Telugu for to reassure, then the Telugu for to satisfy, then say \emph{madhyavarti} again
\end{itemize}

\begin{tcolorbox}[breakable,colback=teal!4,colframe=teal!35!black,title={Before you move on}]
What does \te{మధ్యవర్తి} mean? (\textquotedblleft{}A mediator; a go-between\textquotedblright{}.) Say it once more.
\end{tcolorbox}
